\documentclass[11pt]{article}
\usepackage[T1]{fontenc}
\usepackage[utf8]{inputenc}
\usepackage[margin=1in]{geometry}
\usepackage{lmodern}
\usepackage{microtype}
\usepackage[hidelinks]{hyperref}
\setlength{\parindent}{0pt}
\setlength{\parskip}{6pt}
\setlength{\emergencystretch}{3em}

\begin{document}

\section*{Response to the Academic Editor}

\textbf{Manuscript:} \emph{Proxy-Based Diagnostics of Quantum Oracle Sketching Robustness for Non-IID Sensor and Telemetry Streams}\\
\textbf{Journal:} \emph{Sensors} (MDPI) \textperiodcentered{} \textbf{Manuscript ID:} sensors-4470240

\textbf{Editor's comment.} \emph{``I strongly invite the Authors to not include preprints among references. These works are still undergoing the peer-review process, thus meaning that their content may change prior to publication (if they will be finally published). The Authors can substitute them with related works already published.''}

\textbf{Response.} We thank the Academic Editor for this helpful recommendation. In accordance with the Editor's request, we have completely removed the Zhao et al.\ (2026) preprint, ``Exponential quantum advantage in processing massive classical data'' (arXiv:2604.07639), from the revised manuscript and reference list.

Before editing, we audited the five in-text uses of the preprint and the directly dependent, uncited clauses. We also checked the published works already cited in the manuscript, including Kallaugher (FOCS 2021), Kallaugher, Parekh and Voronova (STOC 2024), Kallaugher, Parekh and Voronova (SOSA 2025), and Huang et al.\ (Science 2022). These papers establish important but different results: problem-specific quantum-streaming space advantages, a black-box quantum-sketch construction, or quantum advantage for learning properties of quantum systems. None establishes the removed preprint's framework-specific classification, dimension-reduction, correlated-data, repetition-overhead, refreshing-time, or classical-hardness claims. We therefore did not mechanically substitute any of them for the removed preprint.

Instead, only the affected statements were revised. Zhao-specific theorem and provenance claims were removed; the remaining peer-reviewed literature is cited only for the general results it actually establishes; and the operational quantities $\tau$, $r$, $T_{\mathrm{eff}}$, $\epsilon$, and the accuracy proxy are now stated unambiguously as constructs of the present paper. The title and the manuscript's ``quantum oracle sketching'' terminology were retained, with the term explicitly defined as the manuscript-level label for this proxy-based diagnostic rather than attributed to a published theorem. The unavoidable Zhao-dependent clauses in the abstract were minimally corrected. Table~1's Zhao-derived ``established theory'' row and the associated memory-reference descriptions were relabelled as paper-defined illustrative scaling references. Figure artwork was not changed; the Figure~6 caption now clarifies the status of the legacy internal labels.

The revised bibliography contains 54 references. No replacement reference was added. No experimental data, methodology, equations, proxy definitions, datasets, seeds, numerical results, figure artwork, or conclusions were changed. All reviewer-requested revisions remain in place, and the complete manuscript compiles cleanly with no unresolved citations or references.

On behalf of all authors,

\textbf{AbdelMoniem Helmy} (corresponding author)\\
\href{mailto:abdelmoniem.hafez@cu.edu.eg}{\nolinkurl{abdelmoniem.hafez@cu.edu.eg}}\\
ORCID 0000-0001-5996-6019

\end{document}
